\section{Hệ Thống Bài Tập Thực Hành Cá Nhân 30 Phút Tại Lớp (Exercise)}

\section{Đề Bài Tập Thực Hành Cá Nhân: Desktop Automation Testing}
\subsection{Môn học: Kiểm Thử Phần Mềm · Chuyên đề: Automation Testing cho Desktop Applications}

- \textbf{Đối tượng:} Sinh viên thực hành cá nhân trực tiếp tại lớp (In-Class Hands-on Lab).
- \textbf{Thời lượng:} Đúng \textbf{30 phút} (5' Chuẩn bị + 12' Bài 1 + 10' Bài 2 + 3' Nộp bài).
- \textbf{Công cụ thực hành:} \textbf{Katalon Studio} (phiên bản Free 11.x) trên hệ điều hành Windows.
- \textbf{Ứng dụng mẫu (AUT):} \texttt{MiniSum.exe} (WinForms C# .NET 8).
- \textbf{Kho tài nguyên chính thức:} [https://github.com/tinphan247/KTPM-2026-Seminar.git](https://github.com/tinphan247/KTPM-2026-Seminar.git)


\subsection{1. Mục Tiêu Bài Thực Hành Cá Nhân}

1. \textbf{Tiếp cận nhanh và độc lập:} Mỗi sinh viên tự thao tác trên máy tính của mình, làm quen với quy trình xây dựng kịch bản kiểm thử tự động cho ứng dụng Windows desktop.
2. \textbf{Sử dụng tài nguyên chuẩn từ GitHub:} Toàn bộ ứng dụng kiểm thử và kho đối tượng (\textit{Object Repository}) được clone trực tiếp từ GitHub, không phụ thuộc vào đường dẫn máy cá nhân của nhóm thuyết trình.
3. \textbf{Nắm vững bản chất Test Oracle & Checkpoint:} Tự tay thiết lập các điểm kiểm tra kết quả kỳ vọng (\texttt{Verify Element Text}) để đánh giá Pass/Fail một cách tự động, chính xác.


\subsection{2. Chuẩn Bị Tài Nguyên (5 Phút Đầu Giờ)}

\subsubsection{Bước 1: Lấy mã nguồn \& Ứng dụng từ GitHub}
Sinh viên mở \textbf{PowerShell} hoặc \textbf{Terminal} trên máy của mình và clone repository:
\begin{lstlisting}
# Clone mã nguồn về thư mục làm việc (ví dụ C:\KTPM hoặc D:\Testing)
git clone https://github.com/tinphan247/KTPM-2026-Seminar.git
cd KTPM-2026-Seminar
\end{lstlisting}

\textit{(Nếu không cài sẵn Git, sinh viên bấm vào link GitHub → chọn \textbf{Code} → \textbf{Download ZIP} và giải nén ra máy).}

\subsubsection{Bước 2: Kiểm tra ứng dụng MiniSum}
- Vào thư mục \texttt{case-study-app/}, nháy đúp chạy thử \texttt{MiniSum.exe} (hoặc \texttt{MiniSum_v3.exe}).
- Cửa sổ đăng nhập \textbf{"MiniSum - Dang nhap"} xuất hiện là máy đã có sẵn \textbf{.NET 8 Desktop Runtime}.
- Tài khoản demo mặc định: \texttt{admin} / Mật khẩu: \texttt{1234}.

\subsubsection{Bước 3: Nạp Starter Kit vào Katalon Studio}
- Mở \textbf{Katalon Studio} $\to$ Mở project mới hoặc nạp sẵn các thư mục tài nguyên từ \texttt{demo-scripts/katalon/} (đã có sẵn \textit{Object Repository} chứa đầy đủ locator của \texttt{txtUser}, \texttt{txtPass}, \texttt{btnLogin}, \texttt{lblMessage}, \texttt{txtA}, \texttt{txtB}, \texttt{btnSum}, \texttt{lblResult}, \texttt{lblError}).


\subsection{3. Nội Dung Chi Tiết Bài Tập Thực Hành (30 Phút)}

\subsubsection{BÀI TẬP 1 (12 Phút): Tự Động Hóa Đăng Nhập \& Checkpoint Thông Báo (4.0 Điểm)}

* \textbf{Tên Test Case:} \texttt{TC_Login_<MSSV>} (Ví dụ: \texttt{TC_Login_23120001}).
* \textbf{Các bước kịch bản:}
  1. \textbf{Khởi động ứng dụng:} Dùng từ khóa \texttt{Start Application}, trỏ đường dẫn tới file \texttt{MiniSum.exe} trong thư mục \texttt{case-study-app/} vừa clone về máy bạn.
  2. \textbf{Ca 1 (Đăng nhập thành công):}
     - Nhập \texttt{txtUser} = \texttt{"admin"}.
     - Nhập \texttt{txtPass} = \texttt{"1234"}.
     - Nhấp chuột nút \texttt{btnLogin}.
     - \textbf{Checkpoint 1 (Xác minh thông báo):} Thêm từ khóa \texttt{Verify Element Text} trên nhãn \texttt{lblMessage} với nội dung kỳ vọng: \texttt{Dang nhap thanh cong}.
  3. \textbf{Ca 2 (Đăng nhập thất bại):}
     - Nhập mật khẩu sai: \texttt{txtPass} = \texttt{"wrongpass"}.
     - Nhấp chuột nút \texttt{btnLogin}.
     - \textbf{Checkpoint 2 (Xác minh báo lỗi):} Kiểm tra \texttt{lblMessage} hiển thị đúng thông báo: \texttt{Sai tai khoan hoac mat khau}.
  4. Chạy kiểm thử và kiểm tra tab \textbf{Log Viewer} báo trạng thái \textbf{PASSED (màu xanh)}.


\subsubsection{BÀI TẬP 2 (10 Phút): Tự Động Hóa Tính Tổng \& Kiểm Thử Giá Trị Biên (6.0 Điểm)}

* \textbf{Tên Test Case:} \texttt{TC_Calculate_<MSSV>}.
* \textbf{Mục tiêu:} Kiểm tra chức năng tính tổng $A + B$ và xác minh các mốc giá trị biên của miền hợp lệ $A, B \in [-99, 99]$.
* \textbf{Các bước kịch bản:}
  1. Sau khi đăng nhập thành công vào màn hình tính toán:
  2. \textbf{Thao tác tính tổng thông thường (2.0 điểm):}
     - Nhập \texttt{txtA} = \texttt{"15"}, nhập \texttt{txtB} = \texttt{"25"}.
     - Nhấp nút \texttt{btnSum}.
     - \textbf{Checkpoint:} Kiểm tra nhãn \texttt{lblResult} có đúng chuỗi: \texttt{Ket qua: 15 + 25 = 40}.
  3. \textbf{Kiểm thử Giá Trị Biên (BVA) (4.0 điểm):}
     - \textbf{Ca biên trên (2.0 điểm):} Nhập $A = 99$, $B = 0 \to$ Bấm \texttt{btnSum} $\to$ Checkpoint xác minh \texttt{lblResult} = \texttt{Ket qua: 99 + 0 = 99}.
     - \textbf{Ca biên dưới (2.0 điểm):} Nhập $A = -99$, $B = 0 \to$ Bấm \texttt{btnSum} $\to$ Checkpoint xác minh \texttt{lblResult} = \texttt{Ket qua: -99 + 0 = -99}.
  4. Thực thi kịch bản và quan sát màn hình tự động điền dữ liệu, tính toán và xác minh.


\subsubsection{PHẦN MỞ RỘNG (Bonus +1.0 Điểm Tùy Chọn)}
* Dành cho sinh viên hoàn thành sớm trong 30 phút:
* Trỏ đường dẫn ứng dụng sang \texttt{case-study-app/MiniSum_v3.exe} (phiên bản có bắt lỗi ngoài biên).
* Nhập giá trị ngoài biên: $A = 100$, $B = 0 \to$ Bấm \texttt{btnSum}.
* \textbf{Thêm Checkpoint:} Dùng \texttt{Verify Element Text} xác minh nhãn lỗi \texttt{lblError} hiển thị thông báo: \texttt{A ngoai mien [-99, 99]}.


\subsection{4. Hướng Dẫn Nộp Bài (3 Phút Cuối Giờ)}

Sinh viên nộp kết quả cho Giảng viên / Trợ giảng lớp học bằng một trong hai cách:
1. \textbf{Cách 1 (Ảnh chụp màn hình Report):} Chụp toàn màn hình Katalon Studio thể hiện cây Test Case mang tên \texttt{<MSSV>}, tab Log Viewer thể hiện các bước thực thi đều \textbf{PASSED}, và cửa sổ ứng dụng đang mở.
2. \textbf{Cách 2 (File Báo cáo HTML):} Mở thư mục \texttt{Reports/} trong project Katalon $\to$ Lấy file \texttt{report.html} (hoặc nén thư mục test case thành \texttt{<MSSV>_LabKatalon.zip}) để gửi qua Google Form / Zalo nộp bài của lớp.


\subsection{5. Barem Chấm Điểm Thực Hành (Thang Điểm 10)}

\begin{table}[H]
\centering
\begin{tabular}{|l|l|l|l|}
\hline
\textbf{STT} & \textbf{Nội dung đánh giá} & \textbf{Tiêu chí chi tiết} & \textbf{Điểm} \\ \hline
**1** & **Môi trường \& Khởi động** & Clone repo thành công, cấu hình đúng đường dẫn `MiniSum.exe` trên máy cá nhân và khởi động được app. & **1.0 đ** \\ \hline
**2** & **Bài 1: Đăng nhập \& Checkpoint** & - Thực hiện ca đăng nhập đúng `admin/1234` + Checkpoint `lblMessage` PASS.<br>- Thực hiện ca đăng nhập sai + Checkpoint báo lỗi PASS. & **3.0 đ**<br>(1.5 đ)<br>(1.5 đ) \\ \hline
**3** & **Bài 2: Tính tổng cơ bản** & Nhập giá trị hợp lệ, bấm tính và Checkpoint `lblResult` hiển thị đúng kết quả. & **2.0 đ** \\ \hline
**4** & **Bài 2: Kiểm thử Biên (BVA)** & - Thực hiện đúng ca biên trên ($A = 99, B = 0$) + Checkpoint PASS.<br>- Thực hiện đúng ca biên dưới ($A = -99, B = 0$) + Checkpoint PASS. & **4.0 đ**<br>(2.0 đ)<br>(2.0 đ) \\ \hline
**5** & **Điểm thưởng (Bonus)** & Bắt được nhãn lỗi `lblError` khi nhập ngoài biên trên `MiniSum\_v3.exe`. & **+1.0 đ**<br>*(tối đa 10đ)* \\ \hline
 & **TỔNG CỘNG** & **Hoàn thành trọn vẹn kịch bản kiểm thử tự động Desktop trong 30 phút** & **10.0 đ** \\ \hline
\end{tabular}
\end{table}



